\section*{Capítulo \ref{ch:b2:enolates-aldol} --- Enóis, enolatos e a reação aldólica}

\begin{solution}{exo:b2:enolates-aldol:1}
Butanona: but-1-en-2-ol \ce{CH2=C(OH)CH2CH3} e but-2-en-2-ol \ce{CH3C(OH)=CHCH3} ((\textit{E}) e
(\textit{Z})). Ciclo-hexanona: ciclo-hex-1-en-1-ol.
\end{solution}

\begin{solution}{exo:b2:enolates-aldol:2}
Etano $<$ propanona $<$ 3-oxobutanoato de etila (10.7) $<$ pentano-2,4-diona (9.0).
\end{solution}

\begin{solution}{exo:b2:enolates-aldol:3}
3-hidroxibutanal, depois but-2-enal sob aquecimento.
\end{solution}

\begin{solution}{exo:b2:enolates-aldol:4}
Propanodioato de dietila, etóxido de sódio, 1-bromobutano; hidrólise ao ácido butilpropanodioico;
o aquecimento remove o \ce{CO2}: \ce{CH3(CH2)3CH2COOH}, ácido hexanoico.
\end{solution}

\begin{solution}{exo:b2:enolates-aldol:5}
Cinético: desprotonação em C6 (o lado \ce{CH2}), com LDA a \qty{-78}{\degreeCelsius}. Termodinâmico:
a ligação dupla mais substituída C1=C2, do lado do carbono que leva a metila, com um alcóxido em seu
álcool à temperatura ambiente.
\end{solution}

\begin{solution}{exo:b2:enolates-aldol:6}
4-fenilbut-3-en-2-ona (benzalacetona), \ce{C6H5CH=CHCOCH3}. O benzaldeído não tem hidrogênio $\alpha$:
ele não pode formar um \omterm{def:b2:enolates-aldol:enolate}{enolato} e é apenas o eletrófilo.
\end{solution}

\begin{solution}{exo:b2:enolates-aldol:7}
O \omterm{def:b2:enolates-aldol:enolate}{enolato} de um grupo metila (C1) ataca a outra carbonila (C5): um anel de cinco membros, a
3-metilciclopent-2-en-1-ona após desidratação. A alternativa, o \omterm{def:b2:enolates-aldol:enolate}{enolato} em C3 atacando C5,
formaria um anel tensionado de três membros.
\end{solution}

\begin{solution}{exo:b2:enolates-aldol:8}
\ce{2CH3COOC2H5 -> CH3COCH2COOC2H5 + C2H5OH} com etóxido de sódio; o produto, desprotonado, é
recuperado como 3-oxobutanoato de etila na acidificação.
\end{solution}

\begin{solution}{exo:b2:enolates-aldol:9}
Propanona, etanal e acetofenona (eles têm um grupo \ce{CH3-CO}; o etanal é \ce{CH3CHO}).
A pentan-3-ona não.
\end{solution}

\begin{solution}{exo:b2:enolates-aldol:10}
O \omterm{def:b2:enolates-aldol:enolate}{enolato} de uma extremidade ataca o éster da outra: 2-oxociclopentano-1-carboxilato de etila, um
anel de cinco membros.
\end{solution}

\begin{solution}{exo:b2:enolates-aldol:11}
Na 3-metilpentan-2-ona o estereocentro (C3) é o carbono $\alpha$: a enolização o torna plano, e a
reprotonação por qualquer das faces o racemiza. Na 4-metil-hexan-2-ona o estereocentro está em C4,
em $\beta$ da carbonila, intocado pela enolização.
\end{solution}

\begin{solution}{exo:b2:enolates-aldol:12}
Propanona com dois equivalentes de benzaldeído em meio básico. Com dois aldeídos diferentes,
condense primeiro a propanona com um aldeído (um equivalente, isolando a benzalacetona) e depois
condense seu grupo metila restante com o segundo. Em um só frasco, os dois aldeídos competem e dão
uma mistura de três dienonas.
\end{solution}

\begin{solution}{pb:b2:enolates-aldol:1}
\textbf{1.} A propanona, seis hidrogênios $\alpha$ (três em cada metila).
\textbf{2.} $\mathrm{CH_3{-}CO{-}CH_2^-} \leftrightarrow \mathrm{CH_3{-}C(O^-){=}CH_2}$.
\textbf{3.} Ele não tem hidrogênio no carbono vizinho de sua carbonila (esse carbono pertence ao anel e
não leva H).
\textbf{4.} O benzaldeído, um aldeído em excesso, é o melhor eletrófilo; o \omterm{def:b2:enolates-aldol:enolate}{enolato} da propanona o
encontra mais vezes que uma carbonila de propanona.
\textbf{5.} Uma cetona tem dois grupos carbonados que empurram densidade eletrônica para seu carbono
carbonílico e dificultam a aproximação; um aldeído tem um só.
\textbf{6.} Basta uma pequena fração de \omterm{def:b2:enolates-aldol:enolate}{enolato}: ele é regenerado à medida que reage.
\textbf{7.} O carbono do \omterm{def:b2:enolates-aldol:enolate}{enolato} se liga ao carbono carbonílico do benzaldeído: o alcóxido da
4-hidroxi-4-fenilbutan-2-ona.
\textbf{8.} Remoção de um hidrogênio $\alpha$ e perda de hidróxido pelo carbono $\beta$: 4-fenilbut-3-en-2-ona.
\textbf{9.} A ligação dupla formada é conjugada ao mesmo tempo com o anel e com a carbonila.
\textbf{10.} O outro grupo metila ainda tem três hidrogênios $\alpha$.
\textbf{11.} \ce{2C6H5CHO + CH3COCH3 -> C6H5CH=CHCOCH=CHC6H5 + 2H2O}.
\textbf{12.} As desidratações, e a precipitação do produto, que deixa a solução.
\textbf{13.} Duas, ambas (\textit{E}).
\textbf{14.} Os isômeros (\textit{Z}) colocam o grupo fenila junto do grupo carbonila: tensão estérica.
\textbf{15.} Duas \ce{C=C}, uma \ce{C=O} e os dois anéis: um sistema conjugado através da molécula.
\textbf{16.} Sua conjugação estendida estreita a diferença $\pi$ (\cref{prop:b2:huckel:gap}): ele absorve na
extremidade azul do visível e parece amarelo.
\textbf{17.} Não: ela é plana (ou quase) com um plano de simetria.
\textbf{18.} 106.12, 58.08 e \qty{234.30}{g/mol}.
\textbf{19.} Benzaldeído $2.1 \times 1.045/106.12 = \qty{20.7}{mmol}$; propanona $0.75 \times
0.7845/58.08 = \qty{10.1}{mmol}$.
\textbf{20.} A propanona precisa de \qty{20.3}{mmol} de benzaldeído, e há \qty{20.7}{mmol} disponíveis:
a propanona é limitante, por pouco.
\textbf{21.} $0.01013 \times 234.30 = \qty{2.37}{g}$.
\textbf{22.} O benzaldeído restante fica no etanol e é eliminado na lavagem.
\textbf{23.} O produto de monocondensação, a benzalacetona.
\textbf{24.} $0.75 \times 2.37 = \boldsymbol{\approx \qty{1.78}{g}}$ de dibenzalacetona.
\end{solution}
